% ID: SEL-D0028-OA-E02
% ESTADO: validada. S029, cotejada con examen y criterios.
\textbf{a)} Igualamos los valores de los tramos en cada empalme.
En $1$, $7=2+a$, de donde $a=5$. En $5$,
$10+a=-25+60+b$, de donde $b=-20$.

\textbf{b)} Con estos parámetros ya hay continuidad.
En $1$ las pendientes laterales son $14$ y $2$, por lo que
$v$ no es derivable allí. En $5$ coinciden, ya que
$2=-2(5)+12$. Por tanto $v'(5)=2$.

Las derivadas son positivas en los interiores de los dos primeros
tramos. En el tercero, $v'(t)=12-2t$ es positiva hasta $6$
y negativa después. La velocidad máxima se alcanza en
$\boxed{t=6}$ y vale $v(6)=16$.
